
On le
munit d'une structure de @@M0@@-@@M1@@-algèbre de la façon suivante : nous avons remarqué plus haut que l'image du foncteur @@M2@@ est dans le sous-groupe des séries formelles en @@M3@@ dont les coefficients en degré non nul sont dans @@M4@@. Utilisant le morphisme canonique : @@M5@@, qui donne : @@M6@@, on en déduit finalement un morphisme :
@@M7@@
qui d'après 1.12 définit un morphisme de foncteurs sur @@M8@@ :
@@M9@@
Les polynômes universels définissant ce morphisme sont les coefficients des @@M10@@~; le @@M11@@-ième coefficient est la composante isobare de poids @@M12@@ de @@M13@@ (démonstration analogue à celle de 1.15), et sera noté @@M14@@.

On étend enfin le morphisme ainsi construit à @@M15@@ de façon naturelle, en utilisant la @@M16@@-structure de @@M17@@. On obtient donc une structure de pré-@@M18@@-anneau sur @@M19@@, définie par les opérations @@M20@@ correspondantes, qui sont telles que :
@@M21@@
et qui seront notées @@M22@@ au lieu de @@M23@@, aucune confusion n'étant possible.

Pour vérifier que c'est une structure de @@M24@@-anneau, on vérifie la commutativité de (3.7.2)~; elle est connue pour les éléments de @@M25@@~; pour les éléments de @@M26@@, elle résulte de celle de (3.7.2) pour @@M27@@ par passage au quotient. On en tire aussitôt le résultat.



% Local macros for Exposé V, §§6--7
\providecommand{\Ghat}{\widehat G}
\providecommand{\Aqhat}{\widehat A_{\mathbb Q}}
\providecommand{\Ch}{\operatorname{Ch}}
\providecommand{\Log}{\operatorname{Log}}
\renewcommand{\id}{\operatorname{id}}
\providecommand{\gr}{\operatorname{gr}}
\providecommand{\Hom}{\operatorname{Hom}}
\providecommand{\epsi}{\varepsilon}
\providecommand{\clgr}{\operatorname{cl}}
\providecommand{\Fil}{\operatorname{Fil}}
\providecommand{\Gr}{\operatorname{Gr}}
\providecommand{\Sym}{\operatorname{Sym}}


